\documentclass[10pt, letterpaper]{article}

% Packages:
\usepackage[
ignoreheadfoot, % set margins without considering header and footer
top=2 cm, % seperation between body and page edge from the top
bottom=2 cm, % seperation between body and page edge from the bottom
left=2 cm, % seperation between body and page edge from the left
right=2 cm, % seperation between body and page edge from the right
footskip=1.0 cm, % seperation between body and footer
% showframe % for debugging 
]{geometry} % for adjusting page geometry
\usepackage{titlesec} % for customizing section titles
\usepackage{tabularx} % for making tables with fixed width columns
\usepackage{array} % tabularx requires this
\usepackage[dvipsnames]{xcolor} % for coloring text
\definecolor{primaryColor}{RGB}{0, 0, 0} % define primary color
\usepackage{enumitem} % for customizing lists
\usepackage{fontawesome5} % for using icons
\usepackage{amsmath} % for math

\usepackage[
pdftitle={Luis Antonio Flores Esquivel CV's},
pdfauthor={Luis Flores},
pdfcreator={Luis Flores},
colorlinks=true,
urlcolor=primaryColor
]{hyperref} % for links, metadata and bookmarks

\usepackage[pscoord]{eso-pic} % for floating text on the page
\usepackage{calc} % for calculating lengths
\usepackage{bookmark} % for bookmarks
\usepackage{lastpage} % for getting the total number of pages
\usepackage{changepage} % for one column entries (adjustwidth environment)
\usepackage{paracol} % for two and three column entries
\usepackage{ifthen} % for conditional statements
\usepackage{needspace} % for avoiding page brake right after the section title
\usepackage{iftex} % check if engine is pdflatex, xetex or luatex

\usepackage{fontawesome5} %for social icons

% Ensure that generate pdf is machine readable/ATS parsable:
\ifPDFTeX
\input{glyphtounicode}
\pdfgentounicode=1
\usepackage[T1]{fontenc}
\usepackage[utf8]{inputenc}
\usepackage{lmodern}
\fi

\usepackage{charter}

% Some settings:
\raggedright
\AtBeginEnvironment{adjustwidth}{\partopsep0pt} % remove space before adjustwidth environment
\pagestyle{empty} % no header or footer
\setcounter{secnumdepth}{0} % no section numbering
\setlength{\parindent}{0pt} % no indentation
\setlength{\topskip}{0pt} % no top skip
\setlength{\columnsep}{0.15cm} % set column seperation
\pagenumbering{gobble} % no page numbering

\titleformat{\section}{\needspace{4\baselineskip}\bfseries\large}{}{0pt}{}[\vspace{1pt}\titlerule]

\titlespacing{\section}{
	% left space:
	-1pt
}{
	% top space:
	0.3 cm
}{
	% bottom space:
	0.2 cm
} % section title spacing

\renewcommand\labelitemi{$\vcenter{\hbox{\small$\bullet$}}$} % custom bullet points
\newenvironment{highlights}{
	\begin{itemize}[
		topsep=0.10 cm,
		parsep=0.10 cm,
		partopsep=0pt,
		itemsep=0pt,
		leftmargin=0 cm + 10pt
		]
	}{
	\end{itemize}
} % new environment for highlights


\newenvironment{highlightsforbulletentries}{
	\begin{itemize}[
		topsep=0.10 cm,
		parsep=0.10 cm,
		partopsep=0pt,
		itemsep=0pt,
		leftmargin=10pt
		]
	}{
	\end{itemize}
} % new environment for highlights for bullet entries

\newenvironment{onecolentry}{
	\begin{adjustwidth}{
			0 cm + 0.00001 cm
		}{
			0 cm + 0.00001 cm
		}
	}{
	\end{adjustwidth}
} % new environment for one column entries

\newenvironment{twocolentry}[2][]{
	\onecolentry
	\def\secondColumn{#2}
	\setcolumnwidth{\fill, 4.5 cm}
	\begin{paracol}{2}
	}{
		\switchcolumn \raggedleft \secondColumn
	\end{paracol}
	\endonecolentry
} % new environment for two column entries

\newenvironment{threecolentry}[3][]{
	\onecolentry
	\def\thirdColumn{#3}
	\setcolumnwidth{, \fill, 4.5 cm}
	\begin{paracol}{3}
		{\raggedright #2} \switchcolumn
	}{
		\switchcolumn \raggedleft \thirdColumn
	\end{paracol}
	\endonecolentry
} % new environment for three column entries

\newenvironment{header}{
	\setlength{\topsep}{0pt}\par\kern\topsep\centering\linespread{1.5}
}{
	\par\kern\topsep
} % new environment for the header

\newcommand{\placelastupdatedtext}{% \placetextbox{<horizontal pos>}{<vertical pos>}{<stuff>}
	\AddToShipoutPictureFG*{% Add <stuff> to current page foreground
		\put(
		\LenToUnit{\paperwidth-2 cm-0 cm+0.05cm},
		\LenToUnit{\paperheight-1.0 cm}
		){\vtop{{\null}\makebox[0pt][c]{
					\small\color{gray}\textit{Last updated in September 2024}\hspace{\widthof{Last updated in September 2024}}
		}}}%
	}%
}%

% save the original href command in a new command:
\let\hrefWithoutArrow\href

% new command for external links:


\begin{document}
	\newcommand{\AND}{\unskip
		\cleaders\copy\ANDbox\hskip\wd\ANDbox
		\ignorespaces
	}
	\newsavebox\ANDbox
	\sbox\ANDbox{$|$}
	
\begin{header}
	\fontsize{25 pt}{25 pt}\selectfont Luis Antonio Flores Esquivel
	
	\vspace{5 pt}
	
	\normalsize
	\faMapMarker* \, Mexico City
	\kern 6 pt \faEnvelope \, \hrefWithoutArrow{mailto:flores.esquivel.luis.antonio2@gmail.com}{flores.esquivel.luis.antonio2@gmail.com}
	\kern 6 pt \faPhone* \, \hrefWithoutArrow{tel:+52-55-51-02-22-44}{+52 5551022244}
	\kern 6 pt \faLinkedin \, \hrefWithoutArrow{https://linkedin.com/in/lflorese}{lflorese}
	\kern 6 pt \faGithub \, \hrefWithoutArrow{https://github.com/LuisFlorEsq}{LuisFlorEsq}
\end{header}

\vspace{-0.3 cm}

	

	
	\section{Education}
	

	
	
	\begin{twocolentry}{
			Graduation date: July 2025
		}
		\textbf{Escuela Superior de Cómputo, Instituto Politécnico Nacional, México}\\
		B.S. in \textit{Artificial Intelligence Engineering}
	
	\end{twocolentry}
	
	\vspace{0.10 cm}
	\begin{onecolentry}
		\begin{highlights}
			\item GPA: 89/100
			\item \textbf{Coursework:} Machine Learning, Deep Learning, Natural Language Processing, Computer Vision
			and Bio-Inspired Computing
		\end{highlights}
	\end{onecolentry}
	
	
	
	
\section{Experience}

\begin{twocolentry}{
		Apr 2025 – Sept 2025
	}
	\textbf{IT Intern – AI Team}, El Puerto de Liverpool S.A.B. de C.V. | Santa Fe, Mexico City
\end{twocolentry}

\vspace{0.10 cm}
\begin{onecolentry}
	\begin{highlights}
		\item Developed and deployed AI-based software solutions using \textbf{FastAPI}, \textbf{Docker}, and \textbf{Google Cloud Platform (Cloud Run, Artifact Registry)}.
		\item Implemented \textbf{computer vision models} (YOLO, DETR-ResNet, segmentation) for automated image editing and object detection.
		\item Built interactive \textbf{user interfaces} with \textbf{Streamlit} and \textbf{Tkinter} to visualize AI outputs and enhance internal tools.
		\item Leveraged \textbf{Google Gen AI} (Gemini, Imagen) for image generation and editing, improving efficiency and content quality.
	\end{highlights}
\end{onecolentry}

	
	
	
	
\section{Major Projects}

\begin{twocolentry}{
		Aug 2024 – June 2025
	}
	\textbf{Facial Recognition–Based Authentication System} | FastAPI, Docker, ArcFace, Azure, Pinecone\\
	\textit{Biometric authentication system for educational institutions to prevent identity fraud during exams and manage access control.}
\end{twocolentry}

\vspace{0.10 cm}
\begin{onecolentry}
	\begin{highlights}
		\item Built an end-to-end facial recognition pipeline using deep learning for feature extraction and identity verification.
		\item Deployed RESTful APIs with FastAPI (Python) integrated into a Kotlin mobile app.
		\item Designed scalable data storage using vector databases for embeddings and relational databases for user records.
	\end{highlights}
\end{onecolentry}

\vspace{0.20 cm}
\begin{twocolentry}{
		Mar 2024 – Jun 2024
	}
	\textbf{DOF – AI (Diario Oficial de la Federación)} | Django, LangChain, Hugging Face, OpenAI API, Pinecone\\
	\textit{Web application using LLMs to summarize and answer queries from public government documents.}
\end{twocolentry}

\vspace{0.10 cm}
\begin{onecolentry}
	\begin{highlights}
		\item Evaluated open-source and proprietary LLMs for question answering and RAG performance.
		\item Developed a conversational chatbot using LangChain and Pinecone for contextual retrieval.
		\item Implemented Django backend services enabling real-time chatbot interaction.
	\end{highlights}
\end{onecolentry}

	
	
	\section{Skills}
	
	
	\begin{onecolentry}
		\textbf{Languages:} Python, Java, C/C++, Matlab
	\end{onecolentry}
	
	\vspace{0.2 cm}
	
	\begin{onecolentry}
		\textbf{Frameworks:} Django, Flask, FastAPI, Langchain
	\end{onecolentry}
	
	\vspace{0.2 cm}
	
	\begin{onecolentry}
		\textbf{Machine Learning and Deep Learning Libraries:} Scikit-learn, Pandas, OpenCV, Numpy, Tensorflow, Pytorch
	\end{onecolentry}
	
	\vspace{0.2 cm}
	
	\begin{onecolentry}
		\textbf{Developer Tools:} Git, Github, Docker, Azure, GCP, LangSmith, Pinecone
	\end{onecolentry}
	
	
	\section{Courses and Certifications}
	
	
\begin{onecolentry}
	\begin{highlights}
		\item \textbf{Google Cloud Computing Foundations Certificate} — Completed coursework covering core \textbf{GCP} services.
		\item \textbf{EVCCCPLN 2023 \& 2024 (Escuela de Verano de Ciencias Cognitivas Computacionales y PLN)} — Participated in introductory workshops on \textbf{NLP}, \textbf{Sentiment Analysis}, \textbf{Semantic Segmentation with CNNs}, and \textbf{Reinforcement Learning}.
	\end{highlights}
\end{onecolentry}

	
	
	
	
\end{document}